\chapter{A Design Space for Diminished Reality under Compositing Constraints}

\section{Introduction}

Chapter 2 established that diminished reality (DR), the selective removal or attenuation of
real-world visual information \citep{mori2017}, has a substantial conceptual literature
but a thin implementation record on consumer hardware. This chapter addresses the technical half of the dissertation's spine
question: \textit{what degree of diminishment control is actually achievable on a consumer video-see-through
(VST) headset, without operating-system-level access to the passthrough layer?}

The answer is not a single number but a structured design space. The central claim of this chapter is
that on a platform such as the Meta Quest 3, the feasibility of any DR effect is determined almost
entirely by \textbf{which layer of the compositing stack an application is permitted to touch}, and that
these permissions stratify into three discrete tiers with sharply different capabilities and costs.
This dissertation carries two vignette modes into the user study, SignPop and Hard Dark, both to
confirmatory testing. A third, Soft Dark, occupies the midpoint of the attenuation axis but was never
shown to a participant (Chapter 6). These were not
designed as arbitrary aesthetic variants. Each occupies a distinct position in this space, and
together they span it. The taxonomy presented here is Contribution C1, and the modes that instantiate
it are described in implementation detail in Chapter 4.

The chapter proceeds as follows. Section 3.2 states the compositing constraint precisely.
Section 3.3 develops the three-tier access taxonomy, the chapter's centrepiece. Section 3.4 makes
explicit what is \textit{impossible}, and why the absence itself is a finding. Section 3.5 derives three
design axes from the taxonomy and locates the three modes on them. Section 3.6 states the novelty
position against prior art, and Section 3.7 states, before an examiner needs to ask, what does and
does not transfer to optical see-through (OST) hardware.

\section{The Compositing Constraint}

On Quest 3, the reconstructed view of the real world that the user experiences ("system passthrough")
is produced and owned by the operating system. An application \textbf{cannot read, modify, or subtract from
the passthrough layer. It can only composite content on top of it}. The passthrough image is never
exposed to application shaders as a texture: it is composited \textit{behind} application content by the OS
compositor, at a stage of the pipeline that applications do not reach. Consequently, the naive
formulation of subtractive DR, "sample the passthrough pixel, darken it, write it back", is not
merely difficult on this platform. It is architecturally impossible.

This is not an incidental limitation but a deliberate privacy and safety boundary, and it is shared
in stronger form by the other major consumer platform: on Apple Vision Pro, camera pixels are
unavailable to consumer applications entirely (raw camera access requires an enterprise-only
entitlement incompatible with App Store distribution), and the only world-appearance control offered
is a global, non-spatial surroundings dimming effect \citep{apple2024}. On true optical see-through
glasses the constraint is harsher still and physical rather than administrative: an additive display
can only \textit{add} light to the optical path. It can superimpose brightness but cannot remove photons
arriving from the world. Darkening is physically unavailable without auxiliary hardware such as
per-pixel dimming layers \citep{itoh2021}.

It is worth being clear about \textit{why} the boundary exists, because its motivation predicts its
persistence. The passthrough image is a continuous camera view of the user's home, workplace, and
bystanders. Granting arbitrary applications read access to it is a privacy exposure, and granting
write access is a safety exposure, since an application could imperceptibly alter the user's view of
the physical world they are walking through. Both concerns strengthen rather than weaken over time,
so research that assumes future read-modify-write access to system passthrough is betting against the
platform vendors' incentives. This dissertation makes the opposite bet: it treats the constraint as
permanent and asks what can be built \textit{inside} it. That is what gives the resulting design space
its claim to durability.

Subtractive attention guidance, dimming, muting, or blacking out the visual periphery, must
therefore be \textit{synthesized} on consumer VST hardware through whatever compositing rights the platform
does grant. Characterising those rights, and the DR capability each one supports, is the purpose of
the taxonomy that follows.

\figplaceholder{60mm}{Diagram of the Quest 3 compositing stack: OS passthrough layer at the back (inaccessible), application render layer(s) in front, compositor combining them. Annotate the three access tiers with arrows showing what each tier can touch.}{The Quest 3 compositing stack and the three tiers of access.}

\section{The Three-Tier Access Taxonomy}

The Quest 3 grants applications three distinct kinds of access relevant to DR, summarised in
Table 3.1 and elaborated below.

\begin{longtable}{p{0.12\textwidth}p{0.20\textwidth}p{0.20\textwidth}p{0.27\textwidth}p{0.15\textwidth}}
\caption{Diminished-reality capability by compositing-access tier on Quest 3.} \\
\hline
\textbf{Tier} & \textbf{Access mechanism} & \textbf{Achievable DR} & \textbf{Ceiling / cost} & \textbf{Modes located here} \\
\hline
\endfirsthead
\hline
\textbf{Tier} & \textbf{Access mechanism} & \textbf{Achievable DR} & \textbf{Ceiling / cost} & \textbf{Modes located here} \\
\hline
\endhead
\textbf{1: OS passthrough styling} &
\verb|OVRPassthroughLayer| Styling API: brightness / contrast / saturation scalars, 3D colour look-up tables, edge tint &
\textbf{Global colour remapping only.} Whole-layer desaturation, tinting, posterisation &
No blur. No per-pixel or spatial masks. No world-locked regions. Layer can never be read. Surface-projected passthrough (per-surface styling) deprecated at SDK v83 &
Not applicable (bounding case) \\
\hline
\textbf{2: Overlay compositing} &
Application draws alpha-blended geometry in front of the passthrough layer &
Dimming, blackout, occlusion. A shaped \textit{absence} of overlay passes OS-native passthrough through untouched &
Can only attenuate or occlude reality. Cannot recolour, blur, or filter it &
Soft Dark, Hard Dark, GranulatedPeriphery, SpotLift. The focus window itself \\
\hline
\textbf{3: Camera re-render (PCA)} &
Passthrough Camera Access: the raw forward camera feed as a GPU texture, with intrinsics and pose &
Arbitrary per-pixel manipulation: blur, desaturation, salience re-grading, glare compression &
Mono 1280$\times$960 at $\sim$85–90$^{\circ}$ horizontal FOV vs the $\sim$110$^{\circ}$ OS view. Added latency. Binocular-fusion burden (Chapter 4). Sustained on-device processing is thermally bounded (Section 3.3.3) &
SignPop (evaluated, built on the ColorPop pipeline), ColorPop, Blur (superseded), ChromaticCool, ConspicuitySqueeze, OutlinedDark \\
\hline
\end{longtable}

\subsection{Tier 1: the styling ceiling}

Tier 1 deserves precision because it prevents an overclaim this dissertation must not make. It is
\textit{not} true that "nothing can be done" to the system passthrough layer: Meta's Passthrough Styling API
allows an application to apply brightness, contrast and saturation adjustments, full 3D colour
look-up tables, and edge tinting to the passthrough layer as a whole \citep{meta2024a}. A global
desaturation of the entire visual field, one conceivable DR primitive, is therefore achievable at
Tier 1 without touching a camera pixel.

What Tier 1 categorically lacks is \textbf{spatial control}. Every styling operation applies to the whole
layer uniformly. There is no mechanism for per-pixel masks, gradients, world-locked regions, or any
effect that treats one part of the visual field differently from another. The one historical
exception, surface-projected passthrough, allowed passthrough to be projected onto
application-defined geometry with per-surface styling. It was deprecated at SDK v83 and is unavailable
going forward. Attention guidance is \textit{definitionally} spatial: its entire purpose is to treat
the focus region differently from the periphery. So Tier 1 alone cannot implement any mode in this
dissertation. \textbf{The absence of spatial control, not the absence of control, is the finding.}

\subsection{Tier 2: diminishment by occlusion}

Tier 2 is ordinary alpha compositing: the application draws translucent or opaque geometry in front
of the passthrough layer. Because the VST display is an opaque screen whose every pixel the
compositor ultimately controls, an application overlay can darken the apparent world without limit,
up to and including full blackout, something physically impossible on additive OST optics. The
constraint runs the other way: a Tier-2 overlay knows nothing about the pixels behind it. It can
attenuate or occlude reality but cannot recolour, sharpen, blur, or respond to its content.

Two of this dissertation's modes live entirely at Tier 2. \textbf{Soft Dark} composites a uniform black
overlay at up to 75 \% opacity across the periphery. \textbf{Hard Dark} takes the same overlay to full
opacity. Both leave a shaped hole, the world-locked focus window, through which the untouched OS
passthrough remains visible.

That hole is the load-bearing architectural insight of the whole system, and it is worth stating in
its general form: \textbf{at Tier 2, the highest-fidelity rendering of reality is an absence of
rendering.} Inside the focus window nothing is drawn at all (overlay alpha = 0), so the user
receives OS passthrough at its full native quality, full $\sim$110$^{\circ}$ field of view, and zero added
latency, precisely where the user is being asked to direct their attention. The manipulation is
confined to the periphery, where visual acuity is lowest and quality loss matters least. This
inverts the naive design of re-rendering the whole view from the camera feed and accepting
camera-grade quality everywhere, and it is why the system escapes the passthrough-acuity ceiling
exactly where task performance depends on seeing well. Alpha punch-through as a mechanism is
documented platform practice \citep{meta2024b} and shipped commercially in
Immersed's Passthrough Portals \citep{immersed2022}, but in all found prior uses the hole reveals
passthrough \textit{surrounded by virtual content}. Using the hole as the \textit{focus region of a diminished
periphery} is, per the prior-art search of Section 3.6, unprecedented.

\subsection{Tier 3: full pixel control at a price}

Tier 3 is the Passthrough Camera Access (PCA) API \citep{meta2025}, which since early 2025 exposes the
headset's forward RGB camera feed to applications as a GPU texture together with camera intrinsics
and pose. At Tier 3 the application is no longer manipulating a proxy for reality. It holds the
camera pixels themselves and can apply arbitrary per-pixel computation: Gaussian blur, luminance-
weighted desaturation, hue-selective salience re-grading, glare compression, everything the ColorPop mode does (Chapter 4).

The price is paid in four currencies:

\begin{enumerate}
\item \textbf{Resolution and field of view.} The accessible feed is a single mono stream at 1280$\times$960
covering roughly 85–90$^{\circ}$ horizontally, against the OS passthrough's $\sim$110$^{\circ}$ stereo view. A Tier-3
re-render is \textit{categorically} lower-fidelity than what the user would otherwise see.
\item \textbf{Latency.} The application-side pipeline (camera capture $\rightarrow$ texture upload $\rightarrow$ shader $\rightarrow$ display)
adds delay that the OS passthrough path, with its dedicated reprojection, does not exhibit.
\item \textbf{Binocular fusion.} The feed is monocular. Presenting it convincingly to two eyes is
non-trivial, and the naive approach measurably fails (the double-vision failure and its
world-direction sampling solution are detailed in Chapter 4).
\item \textbf{Thermal budget.} Sustained on-device processing of the camera stream is bounded. The only
published work to quantify this class of hardware for PCA-style live compositing is a
simulation-based feasibility study, which projects 720p30 operation with thermal throttling after
5–10 minutes \citep{laghari2025}, a projected, not measured, figure, since no benchmark ran on
physical hardware. Tier-3 modes are, on today's hardware, modes for \textit{sessions}, not workdays.
\end{enumerate}

\subsection{The hybrid: composing Tier 2 and Tier 3}

The tiers are not mutually exclusive, and the most capable configuration this work demonstrates is a
composite: the PCA camera feed, shader-processed, rendered onto a head-centred sphere as the
periphery (Tier 3), with the world-locked focus window left transparent so that native OS passthrough
shows through the hole (Tier 2's absence trick). The user experiences full-fidelity reality inside
the window, seamlessly surrounded by a \textit{re-graded interpretation} of reality outside it. This
exploits the quality/FOV asymmetry between the two layers rather than suffering it. The prior-art
search (Section 3.6) found no publication, product, or open-source project that composites system
passthrough with a live shader-processed camera feed in one view, for any purpose. The implementation
is presented in Chapter 4. Its planned quantitative characterisation, legibility across the two
paths, latency, thermal envelope, is the subject of Chapter 5.

\figplaceholder{60mm}{The hybrid composite, exploded-layer diagram: OS passthrough at the back, the effect-processed camera-feed sphere in front with the focus window cut out, and the fused view the user sees.}{The hybrid composite, shown as exploded layers.}

\section{What Is Impossible, and Why It Matters}

A design space is defined as much by its walls as by its rooms. Three impossibilities structure
everything above:

\begin{itemize}
\item \textbf{No read access, at any tier, to the pixels the user actually sees.} Even Tier 3 supplies a
\textit{different, lower-grade copy} of reality, not the passthrough image itself. Any effect that must
respond to what the user sees (content-adaptive dimming, saliency-weighted attenuation as proposed
by \citealp{lu2012}) can respond only to the camera proxy.
\item \textbf{No spatial modulation of the native layer.} The full-quality view can be revealed or hidden
(Tier 2) but never \textit{partially processed}. A "blur the real passthrough periphery" mode, arguably
the perceptually gentlest DR primitive per the foveated-rendering literature \citep{patney2016},
cannot exist on this platform except via the Tier-3 copy.
\item \textbf{No subtraction on OST.} The entire Tier-2 family is VST-specific (Section 3.7).
\end{itemize}

These walls explain the system's shape. The dark modes exist because attenuation-by-overlay is the
only manipulation available at native quality. ColorPop exists because re-grading requires pixels,
and pixels are only available as the Tier-3 copy. The window-as-absence exists because it is the sole
mechanism that delivers native fidelity \textit{and} spatial selectivity simultaneously. Had the platform
offered read-modify-write access to the passthrough layer, none of this architecture would be
necessary. That counterfactual is precisely what makes the taxonomy transferable: any future
platform can be located in it by asking which of the three walls it removes.

\section{Design Axes and the Placement of the Modes}

Three trade-off axes fall out of the taxonomy, and the three study modes were chosen to span them.

\textbf{Fidelity versus control.} Tier 2 offers native fidelity with a single degree of freedom,
attenuation. Tier 3 offers unlimited control at camera-copy fidelity. Hard Dark and Soft Dark sit
at the fidelity pole: the periphery they show, when they show it, is real passthrough, merely
darkened. SignPop sits at the control pole: its periphery is a synthetic re-grading of the camera
copy, built on the ColorPop pipeline.

\textbf{Suppression versus awareness.} Hard Dark removes the peripheral signal entirely: maximal
distractor suppression, zero residual awareness. Soft Dark attenuates to $\sim$75 \% opacity: peripheral
events remain detectable through the veil. SignPop suppresses \textit{selectively}: the underlying ColorPop
pipeline dims the red/orange/green signal band less than other colours by default, and SignPop's
detection gate lifts that band to full visibility, plus a detection aperture, saliency lift and
darkened surround, once a signal lamp or sign is confirmed (Chapter 4, Section 4.5.5). The user-study
safety hypothesis H2b (Chapter 6) lives on exactly this axis. The axis matters because the literature
puts a real cost at its far end. Users asked what they want from diminishment prefer partial
attenuation over complete removal \citep{cheng2022}, and studies of DR in working contexts find
the focus gain paired with a situational-awareness risk \citep{murph2021,mclaughlin2025}. That is why the family keeps a graded point between untouched vision and full blackout, and
why the evaluation measures peripheral awareness rather than assuming it survives.

\textbf{Deployability versus capability.} Soft Dark and Hard Dark run without the camera pipeline at all,
since their shader path samples no texture. They draw negligible GPU power, require no camera
permission, and face no thermal ceiling, so they are deployable for hour-long sessions today. SignPop
inherits Tier 3's full cost structure.

Together, off $\rightarrow$ Soft Dark $\rightarrow$ Hard Dark forms a clean single-variable intensity ladder (overlay alpha
0 $\rightarrow$ 0.75 $\rightarrow$ 1.0), while SignPop extends the space along the orthogonal control axis. This is the
sense in which three modes "span" the space: two axes, both covered, with the intensity axis sampled
at three points.

These three are the modes the study was designed around, though only two of them reached a
participant, and they are not the whole family. Ten peripheral treatments were
built in total, and the remaining seven fill in the same two axes at points the study had no budget
to test. They also divide along the tier boundary, which is the cleanest evidence that the boundary
is real rather than imposed. Four modes never read a camera pixel and are therefore Tier 2: the two
dark modes, plus GranulatedPeriphery, whose grains come from a noise texture rather than the world,
and SpotLift, whose periphery is a plain dark overlay and whose focus-window brightness lift exists
only in the video path, because OS passthrough cannot be brightened. The other six read the camera
feed and are therefore Tier 3: Blur, ColorPop and the SignPop that extends it, ChromaticCool,
ConspicuitySqueeze, and OutlinedDark, which needs five camera taps per fragment just to find the
edges it draws back in. Chapter 4 gives each of them in implementation detail. The point for the
taxonomy is that the tiers were not fitted to three convenient examples. They were derived from
building the whole family and finding that every member's ceiling was set by the tier it had to run
at.

A fourth, orthogonal dimension deserves naming even though this dissertation treats it as design
rather than as an experimental factor: \textbf{time}. Every mode has temporal behaviour: a formation
transition when the window locks, a fade-out when the head moves fast, a fade-in when it settles
(Chapter 4). The literature review found these transition mechanics to be essentially unstudied
in DR. Prior work evaluates diminishment as a static state, not as something that arrives, yields,
and returns. The system therefore exposes every temporal parameter as tunable, and the user study's
exploratory interview asks specifically about the motion-fade experience. A controlled study of DR
transition dynamics is identified in Chapter 7 as future work seeded by this design space rather than
resolved by it.

\figplaceholder{60mm}{Two-axis plot (fidelity/control $\times$ suppression/awareness) with the three modes placed, deployability annotated as marker size or colour.}{Two-axis plot (fidelity/control $\times$ suppression/awareness) with the three modes placed, deployability annotated.}

\section{Novelty Position}

The claims of this chapter are calibrated against a structured prior-art search conducted on
2026-07-05 across ISMAR, CHI, UIST, IEEE VR, DIS and arXiv (2022–2026), Meta developer documentation,
GitHub (including all 136 forks of the upstream PCA samples repository), XR industry press, and the
visionOS and Varjo API surfaces; the full query protocol is archived in the project record. Three
verdicts matter here:

\begin{enumerate}
\item \textbf{The hybrid composite} (system passthrough revealed through a window in a live shader-processed
camera-feed sphere): \textit{nothing found}. No paper, product, or open-source project, for any purpose.
\item \textbf{Peripheral passthrough manipulation for attention guidance on a standalone consumer HMD}:
\textit{nothing found as an implemented system}. The nearest work is concept-level or differently
situated. \citet{mclaughlin2025} evaluate DR distraction-attenuation entirely inside VR
simulation. The DIS 2026 ADHD co-design study is formative, with no prototype. DiminishAR and the
CHI 2026 Vision Pro study target \textit{specific objects/persons} with occlusion rather than the
periphery. \citet{sutton2022} modulate real-world saliency but on a bench-mounted optical
see-through rig that can only add light.
\item \textbf{A dark peripheral vignette over passthrough}: \textit{near precedent, no direct precedent}. Quest OS
v67's "Theatre View" ships an adjustable peripheral dimming vignette, but explicitly not over
passthrough (immersive mode only). Vision Pro's surroundings dimming is global, not windowed.
\end{enumerate}

Each constituent mechanic, taken alone, has documented precedent that this dissertation cites rather
than claims: camera-feed geometry with GPU effects \citep{coviello2025}, a head-centred
fade sphere revealing passthrough \citep{meta2024c}, alpha
punch-through windows \citep{meta2024b,immersed2022}, OS peripheral dimming in immersive mode \citep{meta2024d}, and saliency modulation of reality on OST hardware \citep{sutton2022}. The defensible
novelty claim is therefore precisely bounded: \textbf{the composite}, native passthrough through a
world-locked focus window in an effect-processed camera-feed periphery, exploiting the quality/FOV
asymmetry between layers, \textbf{applied to peripheral attention guidance on a consumer standalone
headset}, together with the world-direction fusion sampling that makes the mono feed binocularly
viewable (Chapter 4). An alpha-blended dark overlay, by itself, is just compositing, and is not
claimed as novel.

Two adjacent patents \citep{patents} describe claim-level ideas in this territory:
attenuating external stimuli while preserving spatial awareness, and blurred external video feeds
composited into VR. Neither shows evidence of implementation. They are noted for
completeness.

\section{What Transfers to Optical See-Through, and What Does Not}

This section makes the non-transfer statement before an examiner makes it. The Tier-2 dark overlay
family, Soft Dark, Hard Dark, and the darkness component of every mode, works because a VST display
is an opaque screen over which the compositor has total authority. \textbf{It does not transfer to optical
see-through glasses}, whose additive optics cannot remove light from the world. On OST, "draw black"
means "draw nothing". A hypothetical OST implementation of peripheral dimming would require optical
hardware assistance (global or segmented dimming layers, per-pixel occlusion masks), which is exactly
Tier-1-style capability relocated into the optics.

What does transfer is, first, the \textbf{taxonomy itself}. OST platforms can be located in it: their
"Tier 2" is empty for subtraction, and their achievable DR collapses onto additive saliency modulation
of the Sutton et al. kind, plus whatever dimming hardware exists. Second, the \textbf{human-factors
findings} of Part H transfer, because what peripheral diminishment does to attention is a property of
the human visual system, not of the display that implements the diminishment. The engineering is
platform-bound. The psychology is not. That division of transferable from non-transferable is itself
a use of the design space, and it is the frame in which Chapter 7 discusses generalisation.

\section{Summary}

On consumer VST hardware, diminished reality is not one problem but three, stratified by compositing
access: a colour-only, spatially-blind styling tier, an attenuation-and-occlusion overlay tier whose
most powerful move is a shaped absence, and a full-control camera-copy tier that pays in fidelity,
latency, fusion and heat. The three study modes span the space this stratification creates, the
hybrid composite exploits the asymmetry between its tiers, and the walls of the space, no reads, no
spatial modulation of the native layer, no subtraction on OST, are findings in their own right.
Chapter 4 descends from this map into the machine.
